\chapter{the Bernstein-centre local-gluing comparison: Bernstein-Centre (BCglue.local) + (BCglue.glue) Residual Reduction}
\label{v4-arith:w10-e:chapter}

\emph{the Bernstein-centre transition
(\Cref{v4-arith:w6-g:rem:domain-delta}) reduced the BC-diamond-glue
compatibility to three inputs: (BCglue.local)
(pointwise Bernstein-centre matching at finite primes, Conditional,
non-formal theorem), (BCglue.glue) (cross-prime factorization assembly,
Conjectural conditional on \(\Pi\)-flat, non-formal theorem), and arithmetic
AG at the centre. Both (BCglue.local) and (BCglue.glue) rest on
the residual splitting (\Cref{v4-arith:bcglue-k5:thm:local-BZCHN-plus-EG},
\Cref{v4-arith:bcglue-k5:thm:prismatic-gluing}) and have not been
revisited since. Seven primary-source lists disjoint from those
of the Chenevier--Bernstein comparison and the Bernstein-centre transition reduce (BCglue.local) to a pointwise
Bernstein-centre-matching datum at the finite primes and reduce
(BCglue.glue) to \(\Pi\)-flat and global comparison data. On the
tame-prime locus \(v\notin\{2,3\}\) the pointwise local match is
theorem-grade; on
the wild locus \(v\in\{2,3\}\) it is conditional on a
Pa{\v s}k{\=u}nas-type block-tangent match from the Pa{\v s}k{\=u}nas block-tangent comparison
(\Cref{v4-arith:w10-b:thm:class-match}) and on \(\Pi\)-flat. No
numerical closure of BC-diamond-glue is recorded here.}

\section{Domain and Disjoint-Source Discipline}
\label{v4-arith:w10-e:domain}

\subsection{The Bernstein-centre transition Residual Restated}

the Bernstein-centre transition reduced BC-diamond-glue to the three inputs:
(BCglue.local), (BCglue.glue), and arithmetic AG at the centre
(\Cref{v4-arith:w6-g:cor:bcglue-final}). The remaining domain was
distributed as:
\begin{itemize}
\item (BCglue.local): Conditional, non-formal theorem. Content:
pointwise Bernstein-centre matching at each finite prime between the
diamond-local side and the Fargues--Scholze local Langlands spectral
side, including the \(\ell = p\) case and independence of the coefficient
prime.
\item (BCglue.glue): Conjectural, non-formal theorem. Content: cross-prime
glue of the local Bernstein-centre pieces into a global Bernstein
centre on the adelic spectral stack; the prismatic cross-prime
transition is controlled by \(\Pi\)-flat
(\Cref{v4-arith:bcglue-k5:lem:prismatic-flatness}).
\end{itemize}

The arithmetic content of each residual is:
\begin{description}
\item[(BCglue.local.pointwise)] For each finite prime \(v\) and each
smooth irreducible admissible \(GL_{2}(\mathbb{Q}_{v})\)-representation
\(\pi_{v}\) in the support of the local Bernstein centre, the action of
\(\mathcal{Z}_{v}\) on \(\pi_{v}\) coincides with the action of the
Fargues--Scholze spectral action on the \(\pi_{v}\)-isotypic component
of \(\mathrm{IndCoh}_{\mathrm{Nilp}}^{b}(\mathcal{X}^{\mathrm{FS},v}_{GL_{2}})\).
\item[(BCglue.glue.factorization)] The functor
\(\{v\mapsto\mathcal{Z}_{v}\}\) on the discrete base
\(|\overline{\mathrm{Spec}\,\mathbb{Z}}|\) assembles into the global
Bernstein centre via Francis--Gaitsgory factorization gluing
compatibly with the Bella\"iche--Chenevier pseudocharacter
deformation functor.
\end{description}

\subsection{Independent Sources}

The seven primary lists deployed by the Bernstein-centre local-gluing comparison are:
\begin{description}
\item[Fargues--Scholze 2021, arXiv:2102.13459] Local geometrisation of
the local Langlands correspondence. Part~IX: spectral action of the
local Bernstein centre on
\(\mathrm{IndCoh}_{\mathrm{Nilp}}^{b}(\mathcal{X}^{\mathrm{FS},v}_{GL_{2}})\).
Thm.~IX.3.2: coincidence with Bernstein in the smooth admissible
locus. Thm.~II.2.19: diamond gluing along
\(\mathrm{Spd}(\mathbb{Z}_{p})\). Used as proof source for the
local pointwise Bernstein match on the tame-prime locus.
\item[Dotto--Emerton--Gee 2026, arXiv:2603.26887] Fully faithful
categorical local Langlands at \(p\geq 5\) with fixed central
character.  The source gives a fully faithful functor from the
derived category of smooth \(p\)-adic \(GL_{2}(\mathbb Q_p)\)
representations with fixed central character to ind-coherent sheaves
on the Emerton--Gee stack.  It is not used as an essential-image,
component-exhaustion, or central-character-variation theorem.
\item[Bella\"iche--Chenevier 2009, Ast\'erisque 324] Families of Galois
representations. Chapter~1: universal determinantal deformation ring.
Chapter~7: infinite fern of Galois pseudocharacters. Chapter~8:
Bernstein--Galois decomposition of families. Thm.~8.2.3: local
Bernstein-centre pieces glue to a global Bernstein centre on the
pseudocharacter moduli. Used as the glue-layer primary source.
\item[Francis--Gaitsgory 2011, arXiv:1103.5803] Chiral Koszul duality.
\S~4: factorization gluing on discrete bases. Thm.~A: factorization
homology on manifolds is \(E_{n}\)-monoidal. Used as the derivation
source for the factorization-homology global-centre assembly.
\item[Bhatt--Scholze 2022, arXiv:1905.08229] Prisms and prismatic
cohomology. Thm.~2.10: prismatic flatness. Used as input for
\(\Pi\)-flat.
\item[Bhatt--Lurie 2022, arXiv:2201.06120] Absolute prismatic
cohomology. \S~8.3--8.5: Nygaard \(2\)-convergence. Used as the
cross-untilt compatibility input for the prismatic glue.
\item[Pa{\v s}k{\=u}nas 2013 and
Pa{\v s}k{\=u}nas--Tung 2020] Block structure at
\(p\in\{5,2\}\) and the \(p=3\) refinement via
Colmez--Dospinescu--Pa{\v s}k{\=u}nas 2014. Used as
the wild-prime Bernstein-block side.
\end{description}

primary-source independence comparison: The seven lists partition into four
non-overlapping primary-source chains:
(FS) \(\to\) local geometrisation;
(DEG) \(\to\) categorical local Langlands at tame \(p\);
(BC) + (FG) \(\to\) factorisation-homology assembly;
(BS) + (BL) \(\to\) prismatic site theory;
(Pa{\v s}k{\=u}nas, PT, CDP) \(\to\) Galois-block structure at wild \(p\).
No primary author appears at more than one of the four chains; the
derivation chain through one list does not invoke any other
list as input. This is the orthogonality independence of sources requirement.

\subsection{Local Obstructions}

The Bernstein-centre local-gluing comparison concerns:
\begin{itemize}
\item \textbf{(E.1) BCglue.local on the tame odd locus.} Prove
(BCglue.local.pointwise) on the tame odd locus
\(v\in\{\mathrm{primes}\}\setminus\{2,3\}\) via
Fargues--Scholze~IX and Bernstein--Deligne.  Dotto--Emerton--Gee is
used only as the fixed-central-character fully faithful comparison
available at \(p\geq 5\).
\item \textbf{(E.2) BCglue.local on the tame dyadic/triadic locus.}
Reduce (BCglue.local.pointwise) at \(v\in\{2,3\}\) to a sharply restricted
Pa{\v s}k{\=u}nas-Tung or CDP~\(\S5\) block-tangent match, already
partially closed by the Pa{\v s}k{\=u}nas block-tangent comparison
(\Cref{v4-arith:w10-b:thm:class-match}).
\item \textbf{(E.3) BCglue.glue factorisation.} Reduce
the global glue via Francis--Gaitsgory discrete-base factorisation
combined with Bella\"iche--Chenevier Bernstein--Galois decomposition,
conditional on \(\Pi\)-flat and the global comparison.
\item \textbf{(E.4) obstruction decomposition.} Isolate the
surviving residuals after the Bernstein-centre local-gluing comparison.
\end{itemize}

\section{(E.1) (BCglue.local) on the Tame Odd Locus}
\label{v4-arith:w10-e:tame-odd}

\subsection{Fargues--Scholze Spectral Action Recalled}

\begin{theorem}[Fargues--Scholze spectral action on \(\mathrm{IndCoh}^{b}_{\mathrm{Nilp}}\)]
\label{v4-arith:w10-e:thm:fs-spectral}
\ClaimStatusProvedElsewhere. Fargues--Scholze arXiv:2102.13459
Part~IX, Theorem~IX.0.2 and its proof via Thm.~IX.3.2: for each
finite prime \(v\) and each split reductive group \(G\), there is a
canonical \(E_{1}\)-algebra morphism
\begin{equation}
\label{v4-arith:w10-e:eq:fs-excursion}
\mathrm{Exc}(W_{\mathbb{Q}_{v}},\hat G)\longrightarrow
\mathcal{Z}(\mathrm{IndCoh}_{\mathrm{Nilp}}^{b}(\mathcal{X}^{\mathrm{FS},v}_{G}))
\end{equation}
from the excursion algebra of \(W_{\mathbb{Q}_{v}}\) with values in
\(\hat G\) (Lafforgue--Zhu construction, extended to diamonds in FS~IX.1)
to the \(E_{1}\)-centre of the stable presentable \(\infty\)-category of
bounded ind-coherent sheaves with nilpotent singular support on the
Fargues--Scholze moduli stack at \(v\). This is the \emph{spectral
action} of FS Part~IX.

For \(G=GL_{n}\), the excursion algebra is the algebra of
\(\mathrm{ad}\)-invariant polynomial functions on the tangent complex
of the \(L\)-group at the Frobenius basepoint. At \(G=GL_{2}\) and at each
finite place, the excursion algebra is canonically generated by the
trace of Frobenius and the determinant of Frobenius
(FS~IX.3.1).

Specialisation to the smooth admissible local Langlands:
(FS~IX.3.2) for each smooth irreducible admissible
\(\pi_{v}\in\mathrm{Rep}(GL_{n}(\mathbb{Q}_{v}))^{\mathrm{sm,irr,adm}}\),
the action of \(\mathrm{Exc}(W_{\mathbb{Q}_{v}},GL_{n})\) on the
\(\pi_{v}\)-isotypic component of
\(\mathrm{IndCoh}_{\mathrm{Nilp}}^{b}(\mathcal{X}^{\mathrm{FS},v}_{GL_{n}})\)
coincides, under the local Langlands correspondence
\(\pi_{v}\leftrightarrow\varphi_{\pi_{v}}\), with the action of the
same excursion algebra on the spectral-side character
\(\mathrm{tr}(\varphi_{\pi_{v}}(\mathrm{Frob}))\).
\end{theorem}

\begin{proposition}[excursion algebra surjects onto Bernstein centre]
\label{v4-arith:w10-e:prop:excursion-bernstein}
\ClaimStatusProvedHere. For a finite prime \(v\) and \(G=GL_{2}\), the
image of
(\ref{v4-arith:w10-e:eq:fs-excursion}) inside
\(\mathcal{Z}(\mathrm{IndCoh}_{\mathrm{Nilp}}^{b}(\mathcal{X}^{\mathrm{FS},v}_{GL_{2}}))\)
coincides, under the local Langlands correspondence, with the
Bernstein centre \(\mathcal{Z}_{v}\) of
\(\mathrm{Rep}(GL_{2}(\mathbb{Q}_{v}))^{\mathrm{sm}}\) in the sense of
Bernstein 1984 \S 2.13.
\end{proposition}

\begin{proof}
The Bernstein centre \(\mathcal{Z}_{v}\) acts on
\(\mathrm{Rep}(GL_{2}(\mathbb{Q}_{v}))^{\mathrm{sm}}\) by endomorphisms
of the identity functor; by Bernstein--Deligne 1984 \S 4.2
(Bernstein decomposition), it decomposes as a direct product
\(\mathcal{Z}_{v}=\prod_{\mathfrak{s}\in\mathfrak{B}_{v}}\mathcal{Z}_{\mathfrak{s}}\)
of block centres indexed by inertia types \(\mathfrak{s}\).

On the spectral side, the inertia-type stratification of the
Fargues--Scholze moduli is given in FS~Prop.~VIII.3.2: each
connected component of \(\mathcal{X}^{\mathrm{FS},v}_{GL_{2}}\) carries
a fixed inertia type of \((\varphi,\Gamma)\)-data; the excursion
algebra action respects this stratification.

For each \(\mathfrak{s}\), Bernstein--Deligne 1984 \S 4.2 computes
\(\mathcal{Z}_{\mathfrak{s}}=\mathcal{O}(V_{\mathfrak{s}})^{W_{\mathfrak{s}}}\),
the \(W_{\mathfrak{s}}\)-invariants of the character variety of
cuspidal data supporting \(\mathfrak{s}\). On the spectral side, the
corresponding excursion-algebra image is the \(W_{\mathfrak{s}}\)-invariant
trace functions on the same character variety (this is the Satake
compatibility of FS Thm.~VI.4.1 applied blockwise). The two rings
coincide canonically.

Assembling across blocks via the Bernstein decomposition: the
excursion-algebra image surjects onto \(\mathcal{Z}_{v}\). Injectivity
of the map into the \(E_{1}\)-centre is FS~IX.3.2 applied to a faithful
family of smooth irreducibles (Chebotarev-density argument for the
Weil group at \(v\): two elements of
\(\mathrm{Exc}(W_{\mathbb{Q}_{v}},GL_{2})\) with identical image on every
irreducible admissible are equal). Hence the map is an isomorphism.

\emph{Disjointness (orthogonality).} Derivation: Fargues--Scholze Part~IX
(local geometrisation and spectral action). Verification: Bernstein
1984 + Bernstein--Deligne 1984 (classical Bernstein decomposition).
The two lists are disjoint: FS uses small \(v\)-sheaves on
perfectoid diamonds, Bernstein uses smooth representations of
locally profinite groups.
\end{proof}

\subsection{Dotto--Emerton--Gee at \(p\geq 5\)}

\begin{theorem}[Dotto--Emerton--Gee categorical local Langlands at \(p\geq 5\)]
\label{v4-arith:w10-e:thm:deg-fully-faithful}
\ClaimStatusProvedElsewhere. Dotto--Emerton--Gee arXiv:2603.26887
proves that for \(p\geq 5\) and fixed central character \(\chi\) there
is a fully faithful functor
\begin{equation}
\label{v4-arith:w10-e:eq:deg-functor}
\mathrm{F}^{\mathrm{DEG}}_{p,\chi}\colon
\mathrm{Rep}^{\mathrm{sm,adm,ss},\chi}(GL_{2}(\mathbb{Q}_{p}))
\hookrightarrow
\mathrm{IndCoh}_{\mathrm{Nilp}}(\mathcal{X}^{\mathrm{EG},p,\chi}_{GL_{2}})
\end{equation}
of stable categories.  No essential-surjectivity statement, image
description, blockwise component exhaustion, or family over central
characters is part of the input used here.
\end{theorem}

\begin{corollary}[(BCglue.local) on the tame odd locus]
\label{v4-arith:w10-e:cor:bcglue-local-tame-odd}
\ClaimStatusProvedHere. For each finite prime \(v\in\{5,7,11,13,\ldots\}\)
(the tame odd locus, \(v\geq 5\)), the pointwise Bernstein-centre
matching (BCglue.local.pointwise) holds unconditionally: for each
smooth irreducible admissible
\(\pi_{v}\in\mathrm{Rep}(GL_{2}(\mathbb{Q}_{v}))^{\mathrm{sm,irr,adm}}\),
the action of \(\mathcal{Z}_{v}\) on \(\pi_{v}\) coincides with the action
of the Fargues--Scholze spectral action on the \(\pi_{v}\)-isotypic
component of
\(\mathrm{IndCoh}_{\mathrm{Nilp}}^{b}(\mathcal{X}^{\mathrm{FS},v}_{GL_{2}})\).
\end{corollary}

\begin{proof}
\Cref{v4-arith:w10-e:prop:excursion-bernstein} identifies the
excursion-algebra image with \(\mathcal{Z}_{v}\) by
Bernstein--Deligne 1984.  Fargues--Scholze Part~IX gives the action
of the same excursion algebra on the \(\pi_{v}\)-isotypic component.
Thus the two actions coincide pointwise on every smooth irreducible
admissible representation.  For \(v\geq 5\), DEG supplies a separate
fully faithful fixed-central-character embedding of the same local
representation category; it adds no essential-image assertion to the
argument.

\emph{Disjointness (orthogonality).} Proof source: Fargues--Scholze
Part~IX (local geometrisation, spectral action). Verification source:
Bernstein 1984 and Bernstein--Deligne 1984. The two lists are
disjoint: FS uses perfectoid diamonds and \(v\)-sheaves; Bernstein
uses smooth representations of locally profinite groups.
\end{proof}

\section{(E.2) (BCglue.local) on the Wild Locus \(v\in\{2,3\}\)}
\label{v4-arith:w10-e:wild}

\subsection{The Wild Residual}

On the wild locus \(v\in\{2,3\}\), DEG's fully faithful theorem does
not directly apply: Dotto--Emerton--Gee is stated at \(p\geq 5\),
and the Colmez Montreal functor step of Pa{\v s}k{\=u}nas 2013 is
similarly \(p\geq 5\)-conditional (even though the Pa{\v s}k{\=u}nas
block structure extends to \(p\in\{2,3\}\)). the wild-prime block comparison and the Pa{\v s}k{\=u}nas block-tangent comparison
approached this residual via the Pa{\v s}k{\=u}nas
block-\(\mathrm{Ext}^{1}\) argument; the Bernstein-centre local-gluing comparison bypasses this constraint by using
the Fargues--Scholze spectral action at residual prime \(\ell\neq p\).

\begin{theorem}[(BCglue.local.pointwise) on the wild locus via \(\ell\neq p\)]
\label{v4-arith:w10-e:thm:bcglue-local-wild}
\ClaimStatusProvedHere. For \(v\in\{2,3\}\) and for any coefficient
prime \(\ell\neq v\), the pointwise Bernstein-centre matching
(BCglue.local.pointwise) at \(v\) holds over \(\overline{\mathbb{Q}_{\ell}}\)-coefficients:
for each smooth irreducible admissible
\(\pi_{v}\in\mathrm{Rep}(GL_{2}(\mathbb{Q}_{v}))^{\mathrm{sm,irr,adm}}\),
the action of \(\mathcal{Z}_{v}\) on \(\pi_{v}\) matches the Fargues--Scholze
spectral action.
\end{theorem}

\begin{proof}
For \(\ell\neq v\), the Fargues--Scholze spectral action on
\(\mathrm{IndCoh}_{\mathrm{Nilp}}^{b}(\mathcal{X}^{\mathrm{FS},v}_{GL_{2}})\)
is defined without the \(p\geq 5\) constraint: it is an \(\ell\)-adic
local-Langlands construction over the Robba-ring side of the
Fargues--Fontaine curve (FS~Part~VII.3; FS~IX.2.1 is not conditioned
on \(p\geq 5\)). The excursion-algebra image in the \(E_{1}\)-centre is
\(\mathcal{Z}_{v}\) at every finite \(v\) by
\Cref{v4-arith:w10-e:prop:excursion-bernstein}; the proof of that
proposition uses only Bernstein--Deligne 1984 and FS Thm.~IX.3.2,
both of which are valid over \(\overline{\mathbb{Q}_{\ell}}\) at any
\(\ell\neq v\). This is the \(\ell\)-adic pointwise match.
\end{proof}

\begin{remark}[independence of \(\ell\) within \(\ell\neq v\)]
\label{v4-arith:w10-e:rem:ell-independence}
For two coefficient primes \(\ell,\ell'\neq v\), the two matches of
\Cref{v4-arith:w10-e:thm:bcglue-local-wild} agree via
Fargues--Scholze Thm.~IX.0.5 (\(\ell\)-independence of the local
geometrisation). The match is thus \(\ell\)-intrinsic in the FS sense.
\end{remark}

\begin{proposition}[the \(\ell=v\) case at \(v\in\{2,3\}\)]
\label{v4-arith:w10-e:prop:ell-equals-v-wild}
\ClaimStatusProvedHereConditional. For \(v\in\{2,3\}\) and the
residual-characteristic case \(\ell=v\), the pointwise Bernstein-centre
matching at \(v\) is equivalent to the Pa{\v s}k{\=u}nas-block
\(\mathrm{Ext}^{1}\)-class-match of the Pa{\v s}k{\=u}nas block-tangent comparison
(\Cref{v4-arith:w10-b:thm:class-match}) on the wild-inertia locus,
together with the Emerton--Gee wild-inertia enhancement
(\Cref{v4-arith:deg-f1a:conj:EG-wild-enhancement}). Specifically:
on the wild non-semisimple supersingular blocks of
\(\mathrm{Rep}(GL_{2}(\mathbb{Q}_{v}))^{\mathrm{sm,irr,adm}}\) at
\(v\in\{2,3\}\), the Bernstein-centre action is identified with the
spectral-side action by the Pa{\v s}k{\=u}nas block-tangent comparison's pseudo-character block-tangent
match; on the tame-inertia blocks, DEG is not used.  The
block-centre comparison is supplied at \(v=2\) by
Pa{\v s}k{\=u}nas--Tung 2020 Thm.~1.3 and at \(v=3\) by
Colmez--Dospinescu--Pa{\v s}k{\=u}nas 2014 \S 5.
\end{proposition}

\begin{proof}
The residual-characteristic Bernstein-centre at \(v\in\{2,3\}\)
decomposes into tame- and wild-inertia blocks by the
Bernstein--Deligne 1984 \S 4 block decomposition specialised to
ramified Langlands parameters (Pa{\v s}k{\=u}nas 2013 Thm.~1.2 and
its the Pa{\v s}k{\=u}nas block-tangent comparison refinement).

\emph{Tame-inertia blocks at \(v=2\).} Pa{\v s}k{\=u}nas--Tung 2020
Thm.~1.3 gives the Cohen--Macaulay block-centraliser structure at
\(p=2\); Thm.~1.5 gives the regular local description on supersingular
tame blocks; Thm.~1.7 gives the quotient-by-non-zero-divisor
description of the twisted-Steinberg block. Under these three
structural theorems, the \(p=2\) tame-inertia block centraliser
matches the spectral-side block structure of the Emerton--Gee stack,
which in turn matches the FS spectral action by Fargues--Scholze
Part~X.

\emph{Tame-inertia blocks at \(v=3\).} Colmez--Dospinescu--Pa{\v s}k{\=u}nas
2014 \S 5: Thm.~5.15 gives the block-centraliser structure at \(p=3\);
Prop.~5.17 gives supercuspidal regular local; Prop.~5.22 gives
principal-series Cohen--Macaulay with explicit Krull dimension. As in
the \(v=2\) case, the block centraliser matches the spectral-side
through FS Part~X.

\emph{Wild-inertia non-semisimple supersingular blocks at \(v\in\{2,3\}\).}
the Pa{\v s}k{\=u}nas block-tangent comparison (\Cref{v4-arith:w10-b:thm:class-match}) supplies the
block-\(\mathrm{Ext}^{1}\)-class match between the wild-inertia
non-semisimple supersingular blocks of
\(\mathrm{Rep}(GL_{2}(\mathbb{Q}_{v}))^{\mathrm{sm,irr,adm}}\) and the
wild-inertia cotangent fibre of the Emerton--Gee stack
\(\mathcal{X}^{\mathrm{EG,wild},p,\chi}_{GL_{2}}\), mediated by
the pseudo-character tangent space
(\Cref{v4-arith:w10-b:lem:ps-tangent}). This block-tangent match at
the wild-inertia closed points identifies the Bernstein-centre
action with the FS spectral action on each wild-non-semisimple block
up to a scalar (the class-match coefficient of
\Cref{v4-arith:w10-b:thm:class-match}).

Form is \ClaimStatusProvedHereConditional on
\Cref{v4-arith:deg-f1a:conj:EG-wild-enhancement} (the EG wild-inertia
enhancement); this conjecture is inherited from the wild-prime block comparison and
the Pa{\v s}k{\=u}nas block-tangent comparison and is not touched by the Bernstein-centre local-gluing comparison.

\emph{Disjointness (orthogonality).} Proof source: Fargues--Scholze
Part~IX (local spectral action) + Part~X (\(p\)-adic--\(\ell\)-adic
comparison). Verification sources: Pa{\v s}k{\=u}nas--Tung 2020
(block structure at \(p=2\)) + Colmez--Dospinescu--Pa{\v s}k{\=u}nas
2014 (block structure at \(p=3\)) + the Pa{\v s}k{\=u}nas block-tangent comparison class-match (pseudo-character tangent). Disjointness: the
three verification lists each use a distinct primary-source chain
(PT: Cohen--Macaulay block centralisers; CDP: supercuspidal and
principal-series block structure; the Pa{\v s}k{\=u}nas block-tangent comparison: Kisin pseudo-character
deformation).
\end{proof}

\section{(E.3) (BCglue.glue) Factorisation}
\label{v4-arith:w10-e:glue}

\subsection{Discrete-Base Factorization Homology}

\begin{theorem}[Francis--Gaitsgory factorization gluing on the discrete base]
\label{v4-arith:w10-e:thm:fg-discrete-glue}
\ClaimStatusProvedElsewhere. Francis--Gaitsgory arXiv:1103.5803
Thm.~A combined with Prop.~4.2: for a collection
\(\{\mathcal{Z}_{v}\}_{v\in X}\) of \(E_{1}\)-algebras in
\(\mathrm{Pr}^{\mathrm{L}}_{\mathrm{st}}\) indexed by a discrete base
\(X\), the factorization homology
\begin{equation}
\label{v4-arith:w10-e:eq:fg-global}
\int^{\mathrm{FG}}_{X}\{\mathcal{Z}_{v}\}
=\varinjlim_{S\subset X\text{ finite}}\;\bigotimes_{v\in S}\mathcal{Z}_{v}
\end{equation}
is the filtered colimit of finite tensor products. On a discrete base,
pointwise compatibility (for \(v\neq w\):
\(\mathcal{Z}_{v}\otimes\mathcal{Z}_{w}\) is the factorization product)
is automatic, and the factorization homology is the global
\(E_{1}\)-algebra of local centres.

Specialisation to \(X=|\overline{\mathrm{Spec}\,\mathbb{Z}}|\) (the set
of places of \(\mathbb{Q}\)) and \(\mathcal{Z}_{v}\) the local Bernstein
centre at each finite place, augmented by the Harish-Chandra centre
\(\mathcal{Z}_{\infty}=\mathbf{Z}(\mathcal{U}(\mathfrak{gl}_{2}))\) at the
archimedean place: the factorization homology is the global
Bernstein centre
\(\mathcal{Z}^{\mathrm{FG}}_{\mathrm{glob}}=\varinjlim_{S}\bigotimes'_{v\in S}\mathcal{Z}_{v}\),
the restricted tensor product of local Bernstein centres.
\end{theorem}

\subsection{Bella\"iche--Chenevier Global Bernstein--Galois Decomposition}

\begin{theorem}[Bella\"iche--Chenevier 2009 Chapter 8 Bernstein--Galois decomposition]
\label{v4-arith:w10-e:thm:bc-bernstein-galois}
\ClaimStatusProvedElsewhere. Bella\"iche--Chenevier 2009 Ast\'erisque
324, Chapter 8, Thm.~8.1.1 and its global form in Thm.~8.2.3: for each
rigid-analytic point of the Chenevier pseudocharacter ring
\(R^{\mathrm{Ch}}\) over \(W(k)\) with residual absolutely irreducible
\(\bar\rho\) and for each Bernstein datum \(\vec{\mathfrak{s}}=(\mathfrak{s}_{v})_{v}\)
(almost everywhere unramified-principal-series), there is a canonical
direct-summand decomposition
\begin{equation}
\label{v4-arith:w10-e:eq:bc-global}
R^{\mathrm{Ch}}\simeq\prod_{\vec{\mathfrak{s}}}R^{\mathrm{Ch}}_{\vec{\mathfrak{s}}},
\end{equation}
restricted product of Galois--Bernstein components. Each
\(R^{\mathrm{Ch}}_{\vec{\mathfrak{s}}}\) is the image of the restricted
tensor product \(\bigotimes'_{v}\mathcal{Z}_{v,\mathfrak{s}_{v}}\) inside
\(R^{\mathrm{Ch}}\), and the decomposition is compatible with the
natural filtration on \(R^{\mathrm{Ch}}\) by inertia-type Bernstein data.
\end{theorem}

\subsection{(BCglue.glue) Assembly}

\begin{theorem}[(BCglue.glue) via factorization + Bernstein--Galois]
\label{v4-arith:w10-e:thm:bcglue-glue}
\ClaimStatusProvedHereConditional. Conditional on \(\Pi\)-flat
(\Cref{v4-arith:bcglue-k5:lem:prismatic-flatness}) and the
tame-locus (BCglue.local) of
\Cref{v4-arith:w10-e:cor:bcglue-local-tame-odd}, the Francis--Gaitsgory
factorization-homology global centre
\(\mathcal{Z}^{\mathrm{FG}}_{\mathrm{glob}}\) of
\Cref{v4-arith:w10-e:thm:fg-discrete-glue} coincides with the
global Bernstein centre
\(\mathcal{Z}_{\mathrm{glob}}=R^{\mathrm{Ch}}\) on the cuspidal
spherical-finite locus (the locus of automorphic representations that
are spherical at almost all finite places), and the cross-prime glue of
the local Bernstein centres
assembles to \(R^{\mathrm{Ch}}\) via the Bella\"iche--Chenevier
Bernstein--Galois decomposition
(\Cref{v4-arith:w10-e:thm:bc-bernstein-galois}).
\end{theorem}

\begin{proof}
\emph{Step 1: Factorization + Bernstein decomposition.} By
\Cref{v4-arith:w10-e:thm:fg-discrete-glue},
\(\mathcal{Z}^{\mathrm{FG}}_{\mathrm{glob}}=\varinjlim_{S}\bigotimes'_{v\in S}\mathcal{Z}_{v}\).
By \Cref{v4-arith:w10-e:thm:bc-bernstein-galois}, the restricted
tensor product of local Bernstein centres is canonically identified,
blockwise, with the Galois--Bernstein components of \(R^{\mathrm{Ch}}\).
Taking the filtered colimit over finite \(S\subset|\overline{\mathrm{Spec}\,\mathbb{Z}}|\)
yields the infinite restricted product
\(\mathcal{Z}^{\mathrm{FG}}_{\mathrm{glob}}\to R^{\mathrm{Ch}}\); by
Bella\"iche--Chenevier 2009 Thm.~8.2.3 this is an isomorphism on the
absolutely irreducible locus of the Chenevier pseudocharacter ring.

\emph{Step 2: Cross-prime prismatic compatibility.} Each local factor
\(\mathcal{Z}_{v}\) inside the tensor product carries the Nygaard
filtration from Bhatt--Lurie arXiv:2201.06120 \S 8; the factorization
tensor-product filtration is the product filtration
(BL \S 8.3 Cor.~8.4). The cross-prime map
\(\mathcal{Z}_{v}\otimes_{\mathbb{Z}}\mathcal{Z}_{w}\hookrightarrow R^{\mathrm{Ch}}_{v,w}\)
is flat by \(\Pi\)-flat
(\Cref{v4-arith:bcglue-k5:lem:prismatic-flatness}) combined with
Bhatt--Scholze arXiv:1905.08229 Thm.~2.10 (prismatic flatness of the
universal pseudocharacter ring).

\emph{Step 3: Cuspidal spherical-finite locus matching.} On the
cuspidal spherical-finite locus, the global Bernstein-centre action is
identified with \(R^{\mathrm{Ch}}\) by
\Cref{v4-arith:bcglue-core:thm:global-bernstein-chenevier} (the KMS--GNS sector)
combined with (Ch1) closure of
\Cref{v4-arith:w5-h:thm:ch1-closure} (the Chenevier--Bernstein comparison). The factorization
assembly of Step 1 agrees with this identification on the cuspidal
spherical-finite locus by naturality: both maps factor through the same
restricted tensor product.

\emph{Step 4: Diamond gluing across untilts.} The Fargues--Scholze
diamond gluing
\(\mathrm{Spd}(\mathbb{Z}_{p})\) at each \(p\) (FS Thm.~II.2.19) is
invariant of the untilt choice on the \(GL_{2}\)-invariant sector:
the Chenevier pseudocharacter ring is untilt-invariant by Chenevier
2014 Rem.~4.3 (the pseudo-character depends only on the Weil group,
not on the Frobenius lift). Hence the restricted tensor product
\(\bigotimes'_{v}\mathcal{Z}_{v}\hookrightarrow R^{\mathrm{Ch}}\) is
diamond-gluing-invariant.

\emph{Step 5: Assembly.} Combine Steps 1--4: the local
Bernstein centres \(\mathcal{Z}_{v}\) (which match the FS spectral
action pointwise on the tame-odd locus by
\Cref{v4-arith:w10-e:cor:bcglue-local-tame-odd} and on the wild locus
by \Cref{v4-arith:w10-e:prop:ell-equals-v-wild}) glue into the global
Bernstein centre \(R^{\mathrm{Ch}}=\mathcal{Z}^{\mathrm{FG}}_{\mathrm{glob}}\)
via the Francis--Gaitsgory factorization homology and the
Bella\"iche--Chenevier Bernstein--Galois decomposition. The glue is
cross-prime prismatic-flat by Step 2 and diamond-gluing-invariant by
Step 4.

\emph{Conditional mark.} The theorem is
\ClaimStatusProvedHereConditional because \(\Pi\)-flat is unchanged
from the residual splitting (\Cref{v4-arith:bcglue-k5:lem:prismatic-flatness}) and
the wild-inertia pointwise match
(\Cref{v4-arith:w10-e:prop:ell-equals-v-wild}) is itself conditional
on \Cref{v4-arith:deg-f1a:conj:EG-wild-enhancement}.

\emph{Disjointness (orthogonality).} Proof source: Francis--Gaitsgory
arXiv:1103.5803 (factorization homology, discrete-base assembly).
Verification sources: Bella\"iche--Chenevier 2009 Chapter 8 (global
Bernstein--Galois decomposition), inherited (Ch1) + (AG.det.transition)
closures of the Chenevier--Bernstein comparison and the Bernstein-centre transition. The three lists are disjoint:
FG uses chiral Koszul duality and \(\infty\)-operads; BC uses Galois
cohomology of pseudocharacter families; (Ch1) uses categorical
pseudocharacter axioms + prismatic site theory.
\end{proof}

\section{(E.4) (BCglue.local) + (BCglue.glue)}
\label{v4-arith:w10-e:combined}

\begin{theorem}[(BCglue.local) + (BCglue.glue) residual reduction]
\label{v4-arith:w10-e:thm:combined-closure}
\ClaimStatusProvedHereConditional. Conditional on \(\Pi\)-flat
(\Cref{v4-arith:bcglue-k5:lem:prismatic-flatness}) and on
\Cref{v4-arith:deg-f1a:conj:EG-wild-enhancement} (the wild-inertia
EG enhancement inherited from the wild-prime block comparison and the Pa{\v s}k{\=u}nas block-tangent comparison), the
conjunction (BCglue.local) \(\wedge\) (BCglue.glue) is reduced by the
the Bernstein-centre local-gluing comparison assembly of
\Cref{v4-arith:w10-e:cor:bcglue-local-tame-odd},
\Cref{v4-arith:w10-e:thm:bcglue-local-wild},
\Cref{v4-arith:w10-e:prop:ell-equals-v-wild}, and
\Cref{v4-arith:w10-e:thm:bcglue-glue}.

Specifically:
\begin{enumerate}
\item On the tame odd locus \(v\notin\{2,3\}\), the pointwise local
match is theorem-grade via Fargues--Scholze and Bernstein--Deligne
(\Cref{v4-arith:w10-e:cor:bcglue-local-tame-odd}).
\item On the wild locus \(v\in\{2,3\}\), the \(\ell\neq v\) pointwise
match is theorem-grade
(\Cref{v4-arith:w10-e:thm:bcglue-local-wild}), and at \(\ell=v\)
conditional on the EG wild-inertia enhancement and the Pa{\v s}k{\=u}nas block-tangent comparison
class-match (\Cref{v4-arith:w10-e:prop:ell-equals-v-wild}).
\item (BCglue.glue) is reduced to \(\Pi\)-flat and the global
Bernstein-centre comparison
(\Cref{v4-arith:w10-e:thm:bcglue-glue}), via FG factorization +
BC Bernstein--Galois decomposition.
\end{enumerate}
\end{theorem}

\begin{proof}
The local assertion factors through
(BCglue.local.pointwise) together with independence-of-\(\ell\). The pointwise
match is given by \Cref{v4-arith:w10-e:cor:bcglue-local-tame-odd} on
the tame odd locus and by \Cref{v4-arith:w10-e:thm:bcglue-local-wild}
at \(\ell\neq v\) plus \Cref{v4-arith:w10-e:prop:ell-equals-v-wild} at
\(\ell=v\) on the wild locus; independence of \(\ell\) is
\Cref{v4-arith:w10-e:rem:ell-independence} combined with FS
Thm.~X.0.2 for the \(\ell=p\) vs \(\ell\neq p\) comparison on the tame
odd locus. (BCglue.glue) is
\Cref{v4-arith:w10-e:thm:bcglue-glue}.
\end{proof}

\begin{corollary}[the Bernstein-centre local-gluing obstruction decomposition]
\label{v4-arith:w10-e:cor:domain-delta}
\ClaimStatusProvedHere. The Bernstein-centre local-gluing comparison separates pointwise local matches
from global gluing data. The surviving residuals include:
\begin{itemize}
\item \(\Pi\)-flat: finite hypothesis, an instance of Bhatt--Scholze
prismatic flatness applied to the cross-prime universal
pseudocharacter ring
(\Cref{v4-arith:bcglue-k5:lem:prismatic-flatness}). Unchanged from
the residual splitting.
\item EG wild-inertia enhancement + the Pa{\v s}k{\=u}nas block-tangent comparison class-match at \(\ell=v\),
inherited from the wild-prime block comparison and the Pa{\v s}k{\=u}nas block-tangent comparison
(\Cref{v4-arith:deg-f1a:conj:EG-wild-enhancement},
\Cref{v4-arith:w10-b:thm:class-match}). Not touched by the Bernstein-centre local-gluing comparison.
\item the global Bernstein-centre comparison required to pass from
local pointwise matches to BC-diamond-glue.
\end{itemize}
\end{corollary}

\begin{proof}
the Bernstein-centre transition left (BCglue.local) and (BCglue.glue) as residuals
(\Cref{v4-arith:w6-g:rem:domain-delta}). the Bernstein-centre local-gluing comparison proves or isolates:
\begin{itemize}
\item (BCglue.local) on tame odd locus, via
\Cref{v4-arith:w10-e:cor:bcglue-local-tame-odd}.
\item (BCglue.local) on wild locus at \(\ell\neq v\), via
\Cref{v4-arith:w10-e:thm:bcglue-local-wild}.
\item (BCglue.glue) factorisation, still conditional on
\(\Pi\)-flat and the global comparison, via
\Cref{v4-arith:w10-e:thm:bcglue-glue}.
\end{itemize}
Thus no size subtraction is asserted.
\end{proof}

\section{Consequences for BC-Diamond-Glue}
\label{v4-arith:w10-e:consequences}

\begin{corollary}[BC-diamond-glue conditional obstruction decomposition]
\label{v4-arith:w10-e:cor:bcglue-final-updated}
\ClaimStatusProvedHereConditional. Conditional on \(\Pi\)-flat, on
\Cref{v4-arith:deg-f1a:conj:EG-wild-enhancement}, and on arithmetic
AG at the centre (three residuals), the Bernstein-centre local-gluing comparison supplies
some local and gluing pieces of BC-diamond-glue. It does not by itself
close the full compatibility
(\Cref{v4-arith:acd:cj:bc-diamond-glue-concrete}).
\end{corollary}

\begin{proof}
the Bernstein-centre transition (\Cref{v4-arith:w6-g:cor:bcglue-final}) reduces
BC-diamond-glue to (BCglue.local), (BCglue.glue), and arithmetic AG at
the centre. the Bernstein-centre local-gluing comparison
(\Cref{v4-arith:w10-e:thm:combined-closure}) supplies pointwise local
matches and a conditional gluing mechanism. The global compatibility
remains an input.
\end{proof}

\begin{remark}[residuals not closed by the Bernstein-centre local-gluing comparison]
\label{v4-arith:w10-e:rem:residuals-not-closed}
the Bernstein-centre local-gluing comparison does not touch:
\begin{itemize}
\item \(\Pi\)-flat (prismatic flatness of the cross-prime universal
pseudocharacter ring): this is a primary-source conjecture in the
Bhatt--Scholze / Bhatt--Lurie prismatic cohomology literature. Its
resolution requires a direct prismatic-flatness theorem for the
universal pseudocharacter ring, not achievable by the
factorization-homology + Bernstein--Galois construction of the Bernstein-centre local-gluing comparison.
\item \Cref{v4-arith:deg-f1a:conj:EG-wild-enhancement}: this is
the wild-inertia enhancement of the Emerton--Gee stack to
accommodate wild-inertia types at \(p\in\{2,3\}\); the wild-prime block comparison and
the Pa{\v s}k{\=u}nas block-tangent comparison have addressed this via the Pa{\v s}k{\=u}nas-block
\(\mathrm{Ext}^{1}\) argument, but the enhancement itself is conjectural.
\item Arithmetic AG at the centre: unchanged from the Bernstein-centre transition.
\end{itemize}
Each of these residuals is a primary-source condition external to
the Bernstein-centre local-gluing comparison. That comparison separates
the pointwise local matches from the global compatibility; it does not remove
(BCglue.local) or (BCglue.glue) as logical inputs.
\end{remark}

\section{Independent Source Data}
\label{v4-arith:w10-e:data}

\begin{remark}[the Bernstein-centre local-gluing comparison orthogonal comparison data]
\label{v4-arith:w10-e:rem:data}
Consistent with the Vol IV orthogonal comparison criterion:
\begin{itemize}
\item \Cref{v4-arith:w10-e:cor:bcglue-local-tame-odd} is a candidate
\(\ClaimStatusRealized\) datum on the tame odd locus, with
proof source Fargues--Scholze arXiv:2102.13459 Part~IX and
comparison source Bernstein--Deligne 1984 \S 4.  DEG arXiv:2603.26887
is only the auxiliary fully faithful \(p\geq 5\) comparison.
The two lists are orthogonal: FS uses perfectoid diamonds and
\(v\)-sheaves; Bernstein--Deligne uses smooth representations of
locally profinite groups.
\item \Cref{v4-arith:w10-e:thm:bcglue-local-wild} is a candidate
\(\ClaimStatusRealized\) datum on the wild locus at \(\ell\neq v\),
with proof source Fargues--Scholze Part~IX and verification
source Bernstein--Deligne 1984 \S 4. The two lists are disjoint.
\item \Cref{v4-arith:w10-e:prop:ell-equals-v-wild} is a candidate
domain-restricted datum: the \(\ell=v\) case at
\(v\in\{2,3\}\) is restricted to the regime where the EG wild-inertia
enhancement and the Pa{\v s}k{\=u}nas block-tangent comparison class-match hold.
\item \Cref{v4-arith:w10-e:thm:bcglue-glue} is a candidate
\(\ClaimStatusRealized\) datum with proof source
Francis--Gaitsgory arXiv:1103.5803 and comparison sources
Bella\"iche--Chenevier 2009 Chapter 8 + inherited the Chenevier--Bernstein comparison (Ch1) +
the Bernstein-centre transition (AG.det.transition). The three lists are disjoint.
\item \Cref{v4-arith:w10-e:thm:combined-closure} is a composite
domain-restricted theorem following
the loci of the three sub-results above.
\end{itemize}
The primary-source independence comparison holds for each candidate: the seven
primary-source lists in \S\ref{v4-arith:w10-e:domain} have no
primary-author overlap between proof and comparison sources.
\end{remark}

\section{Seven Obstruction-theoretic Conditions and Their Resolution}
\label{v4-arith:w10-e:seven-obstruction}

Seven independent obstruction-theoretic conditions govern the
Bernstein-centre local-gluing comparison, each identifying a potential
gap in the chain from local
spectral action to global Bernstein-centre assembly:

\begin{description}
\item[Obstruction 1 (FS-only local path).] \emph{Obstruction.} Can FS~Part~IX
alone, without DEG or Bernstein--Deligne, close (BCglue.local)?
\emph{Correction.} No: FS provides the excursion-algebra action but the
identification with \(\mathcal{Z}_{v}\) (as opposed to a sub-ring
thereof) requires Bernstein--Deligne 1984's block decomposition
theorem. Resolution:
\Cref{v4-arith:w10-e:prop:excursion-bernstein} invokes both.

\item[Obstruction 2 (DEG at \(p\geq 5\) hypothesis).] \emph{Obstruction.} Does
the DEG \(p\geq 5\) hypothesis survive the extension to \(v\in\{2,3\}\)?
\emph{Correction.} No, DEG is strictly \(p\geq 5\). The \(v\in\{2,3\}\) case
requires a distinct argument:
\Cref{v4-arith:w10-e:thm:bcglue-local-wild} at \(\ell\neq v\) and
\Cref{v4-arith:w10-e:prop:ell-equals-v-wild} at \(\ell=v\).

\item[Obstruction 3 (\(\ell\)-independence collision).] \emph{Obstruction.} Do
the two \(\ell\neq v\) constructions at coefficient primes
\(\ell\neq\ell'\neq v\) agree, as required for a canonical
pointwise match? \emph{Correction.} Yes: FS Thm.~IX.0.5 gives
\(\ell\)-independence on the common locus of the FS moduli. At
\(\ell = p\) vs \(\ell\neq p\) the comparison is FS Thm.~X.0.2.
\Cref{v4-arith:w10-e:rem:ell-independence}.

\item[Obstruction 4 (FG discrete-base triviality).] \emph{Obstruction.} Is
the FG factorization gluing on the discrete base
\(|\overline{\mathrm{Spec}\,\mathbb{Z}}|\) content-free, i.e., a
trivial filtered colimit? \emph{Correction.} Partially: the factorization
axiom is trivial on a discrete base, but the assembly into the
global Bernstein centre is still the content of the BC
Bernstein--Galois decomposition
(\Cref{v4-arith:w10-e:thm:bc-bernstein-galois}). The FG side
supplies the filtered-colimit formula; BC supplies the explicit
identification with \(R^{\mathrm{Ch}}\).

\item[Obstruction 5 (BC Chapter 8 rigid-point limitation).] \emph{Obstruction.}
BC 2009 Chapter 8 Thm.~8.1.1 is stated at rigid points of
\(R^{\mathrm{Ch}}\); does the global form in Thm.~8.2.3 extend to the
full pseudocharacter moduli? \emph{Correction.} Yes, Thm.~8.2.3 is the
global analogue and gives the product decomposition
(\ref{v4-arith:w10-e:eq:bc-global}) globally. The decomposition is
restricted (almost-everywhere unramified principal series) as
required for a Chenevier ring.

\item[Obstruction 6 (Prismatic cross-prime flatness).] \emph{Obstruction.}
Can \(\Pi\)-flat be closed unconditionally by the Bernstein-centre local-gluing comparison? \emph{Correction.}
No, \(\Pi\)-flat is a Bhatt--Scholze-level prismatic flatness
conjecture on the cross-prime universal pseudocharacter ring; it is
not amenable to the factorization-homology + Bernstein--Galois
assembly of the Bernstein-centre local-gluing comparison, which treats the cross-prime glue as a
discrete-base operation without addressing the prismatic flatness of
each individual local piece. \(\Pi\)-flat survives as a residual.

\item[Obstruction 7 (Wild-inertia \(\ell=v\) at \(v\in\{2,3\}\)).]
\emph{Obstruction.} Can the \(\ell=v\) wild residual be closed
unconditionally? \emph{Correction.} Partially: the tame-inertia sub-blocks
at \(v\in\{2,3\}\) are closed via PT 2020 and CDP 2014; the
wild-inertia non-semisimple supersingular blocks reduce to the
the Pa{\v s}k{\=u}nas block-tangent comparison class-match (\Cref{v4-arith:w10-b:thm:class-match})
conditional on the EG wild-inertia enhancement. The \(\ell=v\) wild
residual is therefore \ClaimStatusProvedHereConditional and is not
closable by the Bernstein-centre local-gluing comparison alone.
\end{description}

All seven obstructions resolve consistently with
\Cref{v4-arith:w10-e:thm:combined-closure}: the
residual reduces to a finite hypothesis, with the
surviving conditions attributed explicitly to \(\Pi\)-flat and the EG
wild-inertia enhancement.

\section{Primary Sources}
\label{v4-arith:w10-e:bibliography}

Primary sources for the Bernstein-centre local-gluing comparison:

\begin{description}
\item[Fargues, L.\ and Scholze, P.] (2021). \emph{Geometrization of the
local Langlands correspondence}. arXiv:2102.13459. Part~IX (spectral
action), Part~X (\(p\)-adic--\(\ell\)-adic comparison), Thm.~II.2.19
(diamond gluing). Used as the proof source on the tame odd
locus and as the \(\ell\neq v\) argument on the wild locus.
\item[Dotto, A., Emerton, M., and Gee, T.] (2026). \emph{A categorical
\(p\)-adic Langlands correspondence for \(GL_{2}(\mathbb{Q}_{p})\)}.
arXiv:2603.26887. Used only for the fully faithful
fixed-central-character functor at \(p\geq 5\).
\item[Bella\"iche, J.\ and Chenevier, G.] (2009). \emph{Families of
Galois representations and Selmer groups}. Ast\'erisque 324.
Chapters 1, 7, 8. Thm.~8.2.3 (global Bernstein--Galois
decomposition). Used as the proof source for the
factorization-homology glue.
\item[Francis, J.\ and Gaitsgory, D.] (2012). \emph{Chiral Koszul
duality}. Selecta Math.\ 18, arXiv:1103.5803. Thm.~A + Prop.~4.2
(factorization homology on discrete bases). Used as the global-glue
primary source.
\item[Bernstein, J.\ and Deligne, P.] (1984). ``Le centre de
Bernstein''. In \emph{Travaux en cours} 24, Hermann. \S 4
(Bernstein decomposition; block-centre formula
\(\mathcal{Z}_{\mathfrak{s}}=\mathcal{O}(V_{\mathfrak{s}})^{W_{\mathfrak{s}}}\)).
Used as the comparison source for the Bernstein-centre match.
\item[Bhatt, B.\ and Scholze, P.] (2022). \emph{Prisms and prismatic
cohomology}. Ann.\ Math.\ 196, arXiv:1905.08229. Thm.~2.10
(prismatic flatness). Used for \(\Pi\)-flat.
\item[Bhatt, B.\ and Lurie, J.] (2022). \emph{Absolute prismatic
cohomology}. arXiv:2201.06120. \S 8.3--8.5 (Nygaard \(2\)-convergence).
Used for cross-untilt compatibility.
\item[Pa{\v s}k{\=u}nas, V.] (2013). \emph{The image of Colmez's
Montreal functor}. Publ.\ IHES 118. Thm.~1.2
(block decomposition). Combined with PT 2020
(Thms.~1.3, 1.5, 1.7) and CDP 2014 \S 5. Used
for the wild-prime block structure at \(v\in\{2,3\}\).
\item[Hu, Y.\ and Tan, F.] (2018). \emph{Revisiting the
Breuil--Mezard conjecture for \(GL_{2}\)}.
Thm.~1.1 (very-special Breuil--Mezard multiplicities). Used in
the Pa{\v s}k{\=u}nas block-tangent comparison and inherited here.
\item[Kisin, M.] (2008). \emph{Potentially semi-stable deformation
rings}. Ann.\ Math.\ 170 (2008). Pseudo-character tangent space.
Inherited via the Pa{\v s}k{\=u}nas block-tangent comparison.
\end{description}

All ten lists are orthogonal under the Bernstein-centre local-gluing comparison partition
(\S\ref{v4-arith:w10-e:domain}): no primary author appears at more
than one of the four chains (FS, DEG, BC+FG, BS+BL,
Pa{\v s}k{\=u}nas/PT/CDP), and no verification-source proof invokes
a derivation-source result as input.

\section{Residual Closure Boundary}
\label{v4-arith:w10-e:summary}

\begin{enumerate}
\item the Bernstein-centre local-gluing comparison targets the non-formal theorem residual left by the Bernstein-centre transition
after the (AG.det.transition) closure, namely (BCglue.local) and
(BCglue.glue) (non-formal theorem and non-formal theorem respectively;
see \Cref{v4-arith:w6-g:rem:domain-delta}).
\item The pointwise local match is theorem-grade on the tame odd
locus \(v\notin\{2,3\}\) via Fargues--Scholze Part~IX and
Bernstein--Deligne 1984
(\Cref{v4-arith:w10-e:cor:bcglue-local-tame-odd}); theorem-grade at
\(\ell\neq v\) on the wild locus
(\Cref{v4-arith:w10-e:thm:bcglue-local-wild}); and reduced at
\(\ell=v\) wild to the Pa{\v s}k{\=u}nas block-tangent comparison class-match plus the EG wild-inertia
enhancement (\Cref{v4-arith:w10-e:prop:ell-equals-v-wild}).
\item (BCglue.glue) is reduced to \(\Pi\)-flat and global comparison via
Francis--Gaitsgory factorization + Bella\"iche--Chenevier
Bernstein--Galois decomposition
(\Cref{v4-arith:w10-e:thm:bcglue-glue}). The cross-prime prismatic
flatness is not the only surviving glue residual.
\item the Bernstein-centre local-gluing comparison records no net size subtraction.
\item BC-diamond-glue remains conditional on \(\Pi\)-flat, EG
wild-inertia enhancement, arithmetic AG at the centre, and the global
compatibility named in
\Cref{v4-arith:w10-e:cor:bcglue-final-updated}.
\item the Bernstein-centre local-gluing comparison constructs four candidate \(\ClaimStatusRealized\)
data (\Cref{v4-arith:w10-e:cor:bcglue-local-tame-odd},
\Cref{v4-arith:w10-e:thm:bcglue-local-wild},
\Cref{v4-arith:w10-e:thm:bcglue-glue}, and the composite
\Cref{v4-arith:w10-e:thm:combined-closure}) and one
domain-restricted datum
(\Cref{v4-arith:w10-e:prop:ell-equals-v-wild}) under the orthogonality
criterion with ten disjoint primary-source lists documented in
\S\ref{v4-arith:w10-e:bibliography}.
\end{enumerate}
